\newpage
\section{Texto}

\begin{poema}[0.49\textwidth]
  {O Pássaro}
  {Augusto Frederico Schmidt}
  {Em \textit{Aurora Lívida} \cite{schmidt1958aurora}, 1958.}
  O pássaro estremece a noite imóvel \\
  Com suas asas. \\
  O pássaro é o primeiro sinal de vida. \\
  É o primeiro som molhando o espaço. \\[\baselineskip]
  Abre-se na superfície escura uma leve ferida, \\
  Por onde começam a surgir luzes virgens, \\
  Que logo se enfloram, \\
  E no espaço, vigorosas e jovens, se atiram. \\
  São luzes frescas, que têm a forma inicial \\
  De lágrimas deslizando na face da sombra, \\
  Mas que nos campos celestes, \\
  Soltas e livres, tomam as formas de corcéis. \\[\baselineskip]
  O pássaro paira sobre o mar recém-nascido \\
  E inaugura o tempo. \\[\baselineskip]
  E frutificará \\
  Na eternidade prometida.
\end{poema}

\begin{figure}[H]
  \centering
  \begin{minipage}{0.46\textwidth}
    \centering
    \includegraphics[
      page=503,
      width=\linewidth
    ]{../../books/Poesia-Completa.pdf}

    \vspace{0.5em}
    \small
    Outra publicação do poema, em \textit{Poesia Completa}
    \cite[p.~501]{schmidt1995poesiacompleta}.
  \end{minipage}
\end{figure}
